\begin{lemma}[Stable change, Cases 1 and 2, with the exact full tail]
Fix explicitly a finite loopless multigraph \(G\) with exact edge type, a
palette of \(k\) colors, a distinguished edge \(e_0\), an input coloring,
an ordered tree-sequence \(T\), a generating coloring \(\varphi_n\), a
positive-rung ETT \(\mathcal E\) on \(T\), a hierarchy base \(\mathcal B\),
and a hierarchy counterexample \(C\).  Assume the properness,
exactly-one-uncolored, complete lower-rung six-pack, fewer-divider
induction, same-rung equal-divider lexicographic minimality, and
\(2n+11\) estimate recorded by a
\noderef{HierarchyMinimumCounterexampleContext} on these explicit data.

Write \(H\) for the good hierarchy of \(C\), \(q\) for its last stored
index, \(T_n\) for the terminal ladder tree,
\[
 A=\noderef{HierarchyStarredStage}(H,q),\qquad
 T_{n,q}=H.\mathit{stage}(q),\qquad
 D_{n,q}=\noderef{HierarchyStoredColorsAt}(H,q),
\]
and put
\[
 \ell=\noderef{HierarchyFullTailLength}(C),\qquad
 B=\noderef{HierarchyFullTailPrefix}(C,\ell-1).
\]
Let \(\alpha\ne\beta\) be colors missing on \(B\).  Let \(z\) certify a
possibly singleton \noderef{TwoColorPathComponent} \(Q\) under
\(\varphi_n\), suppose that \(Q\) is disjoint from the full
\((\alpha,\beta)\)-component through the first vertex of \(T\) missing
\(\alpha\), and put
\[
 \sigma_n=\noderef{TwoColorComponentSwap}
       (G,\varphi_n,\alpha,\beta,z).
\]
Assume exactly the Cases 1/2 alternative:
\[
\begin{split}
 &q>0\ \hbox{ and }\ \alpha\in\overline{\varphi_n}(T_{n,q}),\\
 &\text{or}\\
 &q=0\ \hbox{ and }\bigl(\alpha\in\overline{\varphi_n}(T_n)
   \ \hbox{or}\ 
   [\alpha\in\overline{\varphi_n}(A)
     \text{ and }\alpha,\beta\notin D_n]\bigr).
\end{split}
\]
Finally assume precisely the following full-tail transport data:
\begin{enumerate}
\item for every exact displayed edge position \(t\) with
\(|A|\leq t<|T|\), \(\sigma_n(T.\mathit{edge}(t))=
\varphi_n(T.\mathit{edge}(t))\);
\item for every \(v\in V(B)\),
\(\overline{\sigma_n}(v)=\overline{\varphi_n}(v)\); and
\item for every \(\eta\in D_{n,q}\), the first position at which \(\eta\)
is missing on \(T\), with the terminal fallback convention, is the same
under \(\sigma_n\) and \(\varphi_n\).
\end{enumerate}
No equality of missing-color sets is assumed at the final vertex of \(T\).

Then \(\sigma_n\) is weakly stable at \(A\) and is a mod-coloring of
\(\mathcal E\).  Moreover, either \(A=T_n\), or there are an input coloring
and an actual ETT presentation of \(A\) under \(\sigma_n\) having maximum
property.  When \(q>0\), there are an input coloring, a maximum ETT on
\(A\) under \(\sigma_n\), a hierarchy base, a hierarchy \(K_A\), and a
good-hierarchy witness \(g_A\) such that
\[
 K_A.\mathit{dividerCount}+1=q,
\]
the new base core and closure are exactly those of \(\mathcal B\), and the
\(j\)-th stage of \(K_A\) is exactly the \(j\)-th stage of \(H\) for every
\(0\leq j\leq q\).  At every retained stored index \(0\leq j<q\), the two
\(D_{n,j}\) sets agree, and
\[
 g_A.\Gamma^j_\eta=C.\mathit{good}.\Gamma^j_\eta
 \quad\text{for every }\eta\in D_{n,j}.
\]
Finally, under \(\sigma_n\), \(T_{n,q}\) is
\[
 \left(\bigcup_{\eta\in D_{n,q}}
   C.\mathit{good}.\Gamma^{q-1}_\eta\right)^{-}\text{-closed}.
\]

In addition, there exist an actual
\noderef{CorrespondingExtendedTashkinovTree} certificate on the full
ordered tree \(T\) under \(\sigma_n\), maximum property for its replacement
ETT, a hierarchy base for that replacement, a full hierarchy \(K\), and a
good-hierarchy witness \(g_K\), such that
\[
              \noderef{SameGoodHierarchy}(C.\mathit{good},g_K).
\]
Thus the full hierarchy has the same base core, the same PE closure when
applicable, the same divider count, ordered stages and divider positions,
and, after the canonical index identification, exactly the same
\(D_{n,j}\) and \(\Gamma^j_\eta\) for every stored \(\eta\).
\end{lemma}
\begin{proof}
Let \(H,q,A,T_n,T_{n,q},D_{n,q},\ell,B\) have the meanings in the
statement.  The context fields are definitionally the properness,
one-uncoloredness, lower-rung, fewer-divider, minimality, and large-set
hypotheses in the approved Section 6 front end.  Apply
\noderef{HierarchyCounterexamplePenultimateElementary} to those six
projections and the maximum, positivity, good-hierarchy, and
non-elementarity fields of the supplied
\noderef{HierarchyCounterexample}.  Its conclusion is exactly
\[
       V(B)\text{ is elementary under }\varphi_n.              \tag{1}
\]
This derives (6.6); it is not a field of the context.

We may therefore apply \noderef{StableChangeCasesOneTwo} with (1), the
two missing colors, the path-component and disjointness hypotheses, and
the displayed Cases 1/2 alternative.  Retain all four parts of its
conclusion:
\[
\begin{array}{ll}
\text{(S1)}&\sigma_n\text{ is weakly stable at }A,\\
\text{(S2)}&\sigma_n\text{ is a mod-coloring of }\mathcal E,\\
\text{(S3)}&A=T_n\text{ or }A\text{ has an actual maximum ETT
presentation under }\sigma_n,\\
\text{(S4)}&\text{when }q>0,\text{ the prefix through }A\text{ has the
stated good-hierarchy certificate and final minus-closedness.}
\end{array}                                                     \tag{2}
\]
These are already the first four asserted conclusions.  It remains to
construct the full certificate without assuming any of its components.

\textbf{The exact old route and its transported tail.}
By \noderef{HierarchyCounterexampleTailAtLeastTwo}, applied to the same
context projections, \(\ell\geq2\).  The ordered inclusions in
\noderef{TashkinovHierarchy} give \(|A|\leq|T|\); hence the defining
subtraction and \(\ell\geq2\) give
\[
 |B|=\min\{|A|+\ell-1,|T|\}=|T|-1,\qquad A\subset B\subset T.   \tag{3}
\]
Here both inclusions are exact ordered
\noderef{TreeSequenceInitialSegment}s, since
\noderef{HierarchyFullTailPrefix} is the corresponding
\noderef{TreeSequencePrefix}; thus all individual parallel-edge identities
are retained.

Consider the final certificate in the
\noderef{TashkinovSeries} carried by \(\mathcal E\).  Positivity excludes
the initial constructor of \noderef{TashkinovSeriesStep}.  If the last
kind is PE, that constructor gives a
\(\varphi_n\)-\noderef{TashkinovClosure} from \(T_n\) to the last series
tree.  If it is SE or RE, it gives an exact edge identity \(f_n\), the
equality \(F_n=\{f_n\}\), an exact
\noderef{TreeSequenceAdjoinEdge} seed \(T_n+f_n\), and a
\(\varphi_n\)-closure from that seed to the last series tree.  The
\noderef{ExtendedTashkinovTree} fields place \(T\) as an exact ordered
initial segment of the same last series tree and strictly above \(T_n\).
Expanding \noderef{TashkinovClosure}, and applying
\noderef{TashkinovAugmentationSequenceRestrict} at the length of \(T\),
therefore gives
\[
\begin{cases}
T_n\longrightarrow T,&\text{if the last kind is PE},\\
T_n+f_n\longrightarrow T,&\text{if it is SE or RE},
\end{cases}                                                     \tag{4}
\]
as finite \(\varphi_n\)-\noderef{TashkinovAugmentationSequence}s.
Every step raises the length by one, by
\noderef{TashkinovAugmentation}.  Since \(A\) is an exact intermediate
prefix of \(T\), induction through (4), comparing the current length with
\(|A|\), splits off the exact suffix
\[
                       A\longrightarrow T.                    \tag{5}
\]
In the SE/RE case with \(A=T_n\), the non-TAA seed edge is kept separate:
the same length comparison instead takes the suffix beginning at
\(T_n+f_n\).  This is precisely the seed required later by the SE/RE
alternative of the mod-coloring definition.

We now transport each one-step augmentation in the suffix of (5).
Write it as \(X\to Y\).  The length clause gives
\(|Y|=|X|+1\), and the exact new edge is
\(g=Y.\mathit{edge}(|X|)\).  Since \(A\subseteq X\subset Y\subseteq T\),
the exact-edge identities in
\noderef{TreeSequenceInitialSegment} identify \(g\) with
\(T.\mathit{edge}(|X|)\), where \(|A|\leq |X|<|T|\).  The first transport
hypothesis therefore preserves the value of \(g\).

The old \noderef{TashkinovAugmentation} also supplies a color \(c\) with
\(\varphi_n(g)=c\) and
\(c\in\noderef{TreeMissingColors}(\varphi_n,X)\).  By the definitions of
\noderef{TreeMissingColors} and \noderef{TreeSequenceVertexSet}, the latter
membership supplies an index \(i<|X|+1\) such that
\(v=X.\mathit{vertex}(i)\) is missing \(c\).  If \(Y=T\),
then \(|X|=|T|-1=|B|\), so \(X=B\) by entrywise equality of equal-length
initial segments.  If \(Y\ne T\), the one-step length inequalities give
\(|X|<|T|-1=|B|\), and the common terminal extension shows
\(X\subset B\).  Unfold this exact
\noderef{TreeSequenceInitialSegment}: it sends the displayed index \(i\)
to the same bounded index of \(B\) and identifies
\(B.\mathit{vertex}(i)=X.\mathit{vertex}(i)=v\).  Unfolding
\noderef{TreeSequenceVertexSet} once more therefore proves
\(v\in V(B)\).  The second transport hypothesis gives
\(c\in\overline{\sigma_n}(v)\), hence
\(c\in\noderef{TreeMissingColors}(\sigma_n,X)\).  Length and
initial-segment clauses are coloring-free.  Consequently the same exact
step \(X\to Y\) is a
\(\sigma_n\)-\noderef{TashkinovAugmentation}.  Induction over the finite
suffix proves the transported route
\[
                       A\longrightarrow T                     \tag{6}
\]
under \(\sigma_n\), or the corresponding route from \(T_n+f_n\) when
\(A=T_n\) in the SE/RE case.  Notice that the final vertex of \(T\) was
never used as a missing-color witness: every predecessor lies in \(B\).
This is why no terminal missing-set equality is required.

\textbf{Eligibility from the terminal ladder tree.}
It remains only to connect the appropriate seed to \(A\).  Unfold (S1)
using \noderef{HierarchyWeaklyStable}.  In the PE case it contains the
\noderef{CombinedStableColoring} relation on \(T_n\) and the PE closure;
in the SE/RE case it contains the ordinary stable relation on \(T_n\).
When \(q>0\), it additionally contains
\noderef{TreeColoringInvariant} on \(A\).  These clauses preserve every
missing-color witness in the old prefix route to \(A\) and every displayed
new-edge value on that route.

For \(q=0\) and terminal kind PE, the parallel field of
\noderef{HierarchyBase} supplies an
\noderef{OrderedTreeUnion} certificate for \(A=T_n\vee R_n\), where \(R_n\)
is the PE closure.  Thus every vertex of \(A\) lies in \(T_n\) or \(R_n\),
and every displayed entry after \(T_n\) is an exact entry of \(R_n\).
For a TAA step inside \(A\), choose its old missing-color witness.  The
union certificate puts that vertex in \(T_n\) or \(R_n\), so the
corresponding \noderef{StableColoring} clause preserves its missingness.
Separately, every displayed entry of \(A\) after \(T_n\) is an exact entry
of \(R_n\).  The \noderef{TashkinovClosure} route defining \(R_n\) says
that the old color of such an edge is missing on an earlier \(R_n\)
prefix, hence is missing on all of \(R_n\); the second stable relation
therefore fixes that incident exact edge.  This verifies each transported
augmentation field by field.  For \(q=0\) and terminal kind SE or RE, the
nonparallel base field gives \(A=T_n\), so there is no prefix to transport.
Repeating the preceding one-step argument on the remaining prefix
transports (4) through \(A\).  Concatenating it with (6) yields exactly
\[
\begin{cases}
T_n\longrightarrow T,&\text{under }\sigma_n\text{ in the PE case},\\
T_n+f_n\longrightarrow T,&\text{under }\sigma_n\text{ in the SE/RE case}.
\end{cases}                                                     \tag{7}
\]
The second line retains the exact identity \(f_n\), the singleton equality
\(F_n=\{f_n\}\), and the same ordered adjoin-edge certificate.

The initial-segment and strictness fields of
\noderef{ExtendedTashkinovTree}, together with (7), are precisely one of
the two eligibility alternatives in the universal clause of (S2), after
unfolding \noderef{ModColoring}.  Hence (S2) produces an actual
\noderef{CorrespondingExtendedTashkinovTree} certificate
\(\mathcal C_T\) on the full ordered tree \(T\).  Apply
\noderef{MaximumPropertyTransport} to (S2), \(\mathcal C_T\), and the
maximum field of \(C\).  The replacement ETT of \(\mathcal C_T\) has
maximum property.                                                     \(\tag{8}\)

\textbf{Rebuilding the full hierarchy.}
If the last kind is PE, (S1) supplies the combined-stability premise and
the first line of (7) supplies the augmentation premise of
\noderef{HierarchySplitter}.  Apply that lemma with the context's
lower-rung projection and the maximum field of \(C\).  Choose its
corresponding ETT certificate (and its transported maximum property) as
\(\mathcal C_T\); it supplies a replacement
\noderef{HierarchyBase} using the same core and PE closure and a
\noderef{TashkinovHierarchy} \(K\) with the same stages and divider
positions as \(H\).

If the last kind is SE or RE, use the certificate from (8).
\noderef{CorrespondingExtendedTashkinovTree} gives equality of the rung
number, every positive ladder tree, every raw stored set, exact
connecting-edge set, and extension kind.  Rebuild the base with its core
equal to the replacement terminal ladder tree, transport the support
position along the equality of that tree with \(T_n\), and copy the unused
closure.  Its terminal missing-color field is reflexive under
\(\sigma_n\); its parallel implication is vacuous because the common
terminal kind is not PE; and its nonparallel field is the chosen core
equality.  Copy the divider count, stages, and divider positions of \(H\).
All nestedness, strictness, divider-prefix, and terminal-stage fields are
unchanged ordered-tree statements.  This constructs a full hierarchy
\(K\) on the replacement ETT, again with the same stages and dividers.

In both cases the common corresponding certificate identifies every raw
stored set.  By \noderef{TashkinovStoredColors}, the computed terminal set
\(D_n\) is the union of those raw sets with the colors missing on \(T_n\)
deleted.  Weak stability identifies the missing-color sets at every
vertex of \(T_n\), so the two computed copies of \(D_n\) agree.  Weak
stability and the second transport premise also identify missing-color
sets at every vertex of \(A\), hence at every starred stage.  Here the
last implication is literal, not an implicit conversion of tree inclusion
to set inclusion: if \(S\) is one of the retained starred stages and
\(v\in V(S)\), unfold \noderef{TreeSequenceVertexSet} to write
\(v=S.\mathit{vertex}(i)\).  The nested exact
\noderef{TreeSequenceInitialSegment}s from \(S\) through \(A\), followed
by \(A\subset B\) from (3), carry \(i\) to a bounded displayed index of
\(B\) with the same vertex value.  A second unfolding of
\noderef{TreeSequenceVertexSet} gives \(v\in V(B)\), where the second
transport premise applies.  Expanding
\noderef{HierarchyStoredColorsAt} therefore gives, at every index \(j\),
\[
                         D^K_{n,j}=D^H_{n,j}.                  \tag{9}
\]
For \noderef{HierarchyStageMissingColors}, unfold its zero case as the
terminal ladder tree and its positive cases as the displayed hierarchy
stages.  Each positive stage in question is an exact prefix of \(A\), so
the identical \noderef{TreeSequenceVertexSet} index argument through
\(B\) identifies the two stage-missing sets as well.

Define \(g_K.\Gamma^j_\eta=C.\mathit{good}.\Gamma^j_\eta\), with the
canonical casts supplied by the equal divider counts.  The two-color and
pairwise-disjointness fields of \noderef{GoodHierarchy} transport
immediately using (9).  We next verify the parallel-base field, assuming
that the common terminal kind is PE, and split according to \(q\).

If \(q>0\), use the actual prefix good-hierarchy witness in (S4).  Its
base core and PE closure are exactly the rebuilt full hierarchy's base
core and closure.  Its retained index zero has, by (S4), exactly the
same \noderef{HierarchyStoredColorsAt} set and the same copied Gamma
pairs as \(g_K\).  Consequently its complete
\noderef{GoodHierarchy}.parallelBase assertion rewrites verbatim to the
parallel-base assertion for \(g_K\).  This transports both the failure of
minus-closedness and the cardinality clause; no assertion about all
boundary-edge colors is involved.

Suppose instead that \(q=0\).  By
\noderef{HierarchyStarredStage}, \(A\) is the base core, while the first
stage after it is \(T\) by the terminal-stage field of
\noderef{TashkinovHierarchy}.  The exact strict inclusions in (3) give
\(|A|<|T|\).  The transported route (6) is nonempty; take its first step
\(A\to Y\), remembering both
the old and transported augmentation certificates.  The old
\noderef{TashkinovAugmentation} supplies
\(g=Y.\mathit{edge}(|A|)\), a color \(c\), and
\(\varphi_n(g)=c\) with \(c\in\overline{\varphi_n}(A)\).  The exact
prefix identities identify \(g\) with \(T.\mathit{edge}(|A|)\), so the
first tail-edge transport hypothesis gives \(\sigma_n(g)=\varphi_n(g)=c\).

To transport the missingness of \(c\), expand
\noderef{TreeMissingColors} and \noderef{TreeSequenceVertexSet} and choose
an index \(i<|A|+1\) such that \(c\) is missing at
\(v=A.\mathit{vertex}(i)\).  The exact prefix \(A\subset B\) carries this
same displayed index and vertex into \(B\); unfolding
\noderef{TreeSequenceVertexSet} on \(B\) makes the second tail hypothesis
applicable.  It says that \(c\) is also missing at \(v\) under
\(\sigma_n\), hence \(c\in\overline{\sigma_n}(A)\).  We have therefore
proved
\[
 \varphi_n(g)=\sigma_n(g)=c,qquad
 c\in\overline{\varphi_n}(A)\cap\overline{\sigma_n}(A).       \tag{9a}
\]

We record that \(g\) is an exact boundary identity of \(A\).  Unfold
\noderef{TreeSequence} for \(Y\).  Its attachment field says that \(g\)
joins the newly displayed vertex at position \(|A|+1\) to a vertex at an
earlier displayed position.  The exact prefix \(A\subset Y\) identifies
the latter with a vertex of \(A\), while vertex injectivity and the
one-step length equality show that the new vertex is not in \(A\).
Unfolding \noderef{TreeSequenceVertexSet} and then
\noderef{BoundaryEdges} proves \(g\in\partial_G(V(A))\); (9a) and
\noderef{BoundaryColorEdges} then give
\[
             g\in\partial_{\sigma_n,c}(V(A)).                 \tag{9b}
\]

It remains to show that \(c\) is not reserved at index zero.  Suppose to
the contrary that
\[
 c\in\bigcup_{\eta\in\noderef{HierarchyStoredColorsAt}(H,0)}
                  C.\mathit{good}.\Gamma^0_\eta .              \tag{9c}
\]
Choose the corresponding \(\eta\).  Expanding
\noderef{HierarchyStoredColorsAt}, membership of \(\eta\) says in
particular that \(\eta\notin\overline{\varphi_n}(A)\).  By
\noderef{TreeMissingColors} and \noderef{TreeSequenceVertexSet}, \(\eta\)
is therefore missing at none of the displayed vertices through position
\(|A|\): the exact prefix \(A\subset T\) from (3) identifies each such
vertex of \(T\) with the corresponding displayed vertex of \(A\).  Unfold
\noderef{HierarchyFirstMissingPosition}.  If an
\(\eta\)-missing vertex exists on \(T\), its least position is strictly
after \(|A|\); if none exists, the definition returns the terminal
fallback position \(|T|\), which is also strictly after \(|A|\).

Thus the exact position \(r=|A|\) satisfies
\[
 |A|\le r<\min\{
 \noderef{HierarchyFirstMissingPosition}(\varphi_n,T,\eta),|T|\}.
\]
The index-zero next stage is \(T\), and the old augmentation identifies
its edge there with \(g\) of old color \(c\).  Hence the definition of
\noderef{HierarchySegmentColors} puts
\(c\in\operatorname{Seg}(H,0,\eta)\).  But (9c) puts
\(c\in C.\mathit{good}.\Gamma^0_\eta\), whereas old
\noderef{GoodHierarchy}.avoidance makes that Gamma pair disjoint from
this segment.  This contradiction proves that \(c\) lies outside the
stage-zero Gamma union.

Now unfold \noderef{TreeColorSetMinusClosed} and
\noderef{TreeColorSetClosed}.  The missing-color clause in (9a), the
just-proved exclusion from the reserved union, and the exact colored
boundary witness (9b) show that the core is not minus-closed under
\(\sigma_n\).  For the separate cardinality clause, the two
\noderef{StableColoring} conjuncts of the combined stability in (S1)
identify vertexwise missing-color sets on \(T_n\) and the PE closure.
Together with (9) and the copied Gamma pairs, this identifies the two
finite-set differences in the old and new cardinality-\(4\) assertions.
Thus the old cardinality bound transports.  This uses no claim that all
boundary edges of the union core retain their colors.

Suppose \(q>0\).  Use the prefix good hierarchy in (S4) for all
unique-refresh and minus-closed clauses strictly before the last divider.
Its exact retained stage, stored-set, and Gamma equalities identify these
clauses with those just defined for \(K\).  The unique-refresh clause at
the last divider transports directly from \(C.\mathit{good}\): all vertices
in the relevant stage difference lie in \(A\subset B\), so the second
transport premise identifies their missing sets, while (9) identifies the
quantified stored colors and the copied Gamma pairs.  The remaining last
minus-closed clause is exactly the separately stated minus-closed
certificate in (S4).  When \(q=0\), there are no unique-refresh or
minus-closed fields.  Thus only avoidance remains.

For an index \(j<q\), avoidance is already part of the prefix good hierarchy
in (S4).  Fix instead \(\eta\in D_{n,q}\) and
\(c\in C.\mathit{good}.\Gamma^q_\eta\).  Old avoidance says, after
unfolding \noderef{HierarchySegmentColors}, that \(c\) is stage-missing
and that there is no exact position \(r\) satisfying
\[
 |A|\leq r<\min\{
 \noderef{HierarchyFirstMissingPosition}(\varphi_n,T,\eta),|T|\}
 \quad\text{and}\quad \varphi_n(T.\mathit{edge}(r))=c.         \tag{10}
\]
The third transport premise applies to this particular stored \(\eta\)
and makes the first-position endpoint in (10) identical under the two
colorings.  At every candidate \(r\), the first transport premise makes
the exact displayed edge value identical.  Hence a witness placing \(c\)
in the new segment would, point by point, be a witness placing \(c\) in
the old segment, contradicting (10).  The stage-missing membership
transports by the vertexwise equality already used above.  This proves
the last avoidance clause.  It does not assert a generic equality of
whole hierarchy segment-color sets.

All fields of \(g_K:\noderef{GoodHierarchy}(K)\) have now been supplied.
The construction gives equality of the base core, divider count, stage
and divider arrays, and the PE closure when applicable; (9) and the
definition of \(g_K\) give exact stored-set and Gamma agreement at every
index.  These are precisely the two clauses of
\noderef{SameGoodHierarchy}(C.\mathit{good},g_K).  Together with (2) and
(8), this is the full asserted conclusion.
\end{proof}
